% (c) Nikita Lisitsa, lisyarus@gmail.com, 2026

\documentclass[10pt]{beamer}

\usepackage[T2A]{fontenc}
\usepackage[russian]{babel}
\usepackage{minted}

\usepackage[most]{tcolorbox}
\tcbuselibrary{minted}

\usepackage{graphicx}
\graphicspath{ {./images/} }

\usepackage{adjustbox}

\usepackage{color}
\usepackage{soul}
\usepackage{multicol}
\usepackage{hyperref}
\usepackage{tikz}
\usetikzlibrary{arrows.meta,positioning,calc}

\usetheme{metropolis}

\definecolor{red}{rgb}{1,0,0}
\definecolor{codebg}{RGB}{29,35,49}
\setminted{fontsize=\footnotesize}

\makeatletter
\newcommand{\slideimage}[1]{
  \begin{figure}
    \begin{adjustbox}{width=\textwidth, totalheight=\textheight-2\baselineskip-2\baselineskip,keepaspectratio}
      \includegraphics{#1}
    \end{adjustbox}
  \end{figure}
}
\makeatother

\title{Язык программирования C++}
\subtitle{Лекция 22: динамически загружаемые модули (плагины), библиотека Boost, работа с сетью}
\author{Лисица Никита}
\date{2026}

\setbeamertemplate{footline}[frame number]

\usemintedstyle{solarized-light}

\begin{document}

\frame{\titlepage}

\begin{frame}[fragile]
\frametitle{Библиотеки: напоминание}
\begin{itemize}
\usemintedstyle{tango}
\item При линковке программы можно линковать \textit{статические} и \textit{динамические библиотеки}
\pause
\item Код из \textit{статических} библиотек (в бинарном виде!) \underline{копируется} в наш исполняемый файл
\pause
\item На код из \textit{динамических} библиотек \underline{ссылается} наш исполняемый файл
\pause
\item При запуске исполняемого файла операционная система сама подгружает необходимые динамические библиотеки
\end{itemize}
\end{frame}

\begin{frame}[fragile]
\frametitle{Плагины}
\begin{itemize}
\usemintedstyle{tango}
\item Иногда хочется \underline{в процессе работы} программы (а не при старте!) загрузить \underline{заранее неизвестную} библиотеку
\pause
\item Как правило, это т.н. \textit{плагин} -- дополнительная, заранее не предусмотренная функциональность программы
\pause
\item Некоторые системные библиотеки иногда тоже загружаются этим способом (напр. загрузчики OpenGL/Vulkan)
\pause
\item Как это сделать?
\end{itemize}
\end{frame}

\begin{frame}[fragile]
\frametitle{Плагины: загрузка/выгрузка}
\begin{itemize}
\usemintedstyle{tango}
\item Загрузить файл динамической библиотеки по заданному пути: \mintinline{cpp}|dlopen| (Linux, MacOS) / \mintinline{cpp}|LoadLibraryW| (Windows)
\pause
\item Возвращает некий идентификатор библиотеки
\pause
\item Когда библиотека больше не нужна, её следует отгрузить: \mintinline{cpp}|dlclose| (Linux, MacOS) / \mintinline{cpp}|FreeLibrary| (Windows)
\end{itemize}
\end{frame}

\begin{frame}[fragile]
\frametitle{Плагины: загрузка функций}
\begin{itemize}
\usemintedstyle{tango}
\item Загрузить из динамической библиотеки функцию по её имени: \mintinline{cpp}|dlsym| (Linux, MacOS) / \mintinline{cpp}|GetProcAddress| (Windows)
\pause
\item \textbf{NB:} ожидается имя функции после \textit{name mangling}'а! В C++ коде лучше всего занести экспортируемую функцию в \mintinline{cpp}|extern "C"|
\pause
\item Возвращает \mintinline{cpp}|void *|, ничего не знает о настоящем типе функции!
\pause
\item Чтобы вызвать функцию, придётся явно привести тип к указателю на функцию:
\end{itemize}
\usemintedstyle{lightbulb}
\begin{tcblisting}{listing engine=minted, minted language=cpp, minted options={fontsize=\scriptsize}, listing only, colback=codebg, colframe=codebg, rounded corners}
// Код программы
auto func = static_cast<void(int)>(dlsym(library, "myfunc"));
func(42);
\end{tcblisting}
\end{frame}

\begin{frame}[fragile]
\frametitle{Плагины: загрузка функций}
\begin{itemize}
\usemintedstyle{tango}
\item Неудобно помнить о типах большого числа функций
\pause
\item Если не угадать тип функции, произойдёт UB (вероятнее всего, креш)
\pause
\item Приходится заранее договариваться об \textit{интерфейсе} плагина
\end{itemize}
\end{frame}

\begin{frame}[fragile]
\frametitle{Плагины: интерфейс через классы}
\begin{itemize}
\usemintedstyle{tango}
\item Удобно выделить интерфейс плагина в виде ООП-интерфейса, т.е. класса с виртуальными функциями:
\usemintedstyle{lightbulb}
\begin{tcblisting}{listing engine=minted, minted language=cpp, minted options={fontsize=\scriptsize}, listing only, colback=codebg, colframe=codebg, rounded corners}
// Общий header программы и плагина
class IPlugin {
public:
  void initialize(const ProgramContext&) = 0;
  void apply(Data&) = 0;
  void serialize(const Data&, std::ostream&) = 0;

  std::string name() const = 0;
  std::string describe() const = 0;

  ~IPlugin() = 0;
};
\end{tcblisting}
\end{itemize}
\end{frame}

\begin{frame}[fragile]
\frametitle{Плагины: интерфейс через классы}
\begin{itemize}
\item Тогда в самом плагине достаточно сделать функцию, создающую реализацию интерфейса:
\usemintedstyle{lightbulb}
\begin{tcblisting}{listing engine=minted, minted language=cpp, minted options={fontsize=\scriptsize}, listing only, colback=codebg, colframe=codebg, rounded corners}
// Код плагина
class MyPlugin : public IPlugin { ... };

extern "C" IPlugin* create() {
  return new MyPlugin();
}
\end{tcblisting}
\end{itemize}
\end{frame}

\begin{frame}[fragile]
\frametitle{Плагины: интерфейс через классы}
\begin{itemize}
\item Для загрузки плагина нужно загрузить эту функцию и вызвать:
\usemintedstyle{lightbulb}
\begin{tcblisting}{listing engine=minted, minted language=cpp, minted options={fontsize=\scriptsize}, listing only, colback=codebg, colframe=codebg, rounded corners}
// Код программы
auto create = static_cast<IPlugin*()>(
  dlsym(library, "create"));

IPlugin* plugin = create();
plugin->initialize();

std::cout << "Loaded plugin "
  << plugin->name() << std::endl;
\end{tcblisting}
\end{itemize}
\end{frame}

\begin{frame}[fragile]
\frametitle{Плагины: удаление}
\begin{itemize}
\usemintedstyle{tango}
\item В общем случае, нет гарании, что динамическая память (куча) процесса и динамической библиотеки -- это одно и то же
\pause
\item \(\Longrightarrow\) Объекты, выделенные через \mintinline{cpp}|malloc/new| в динамической библиотеке, нужно удалять через \mintinline{cpp}|free/delete| в коде библиотеки
\usemintedstyle{lightbulb}
\begin{tcblisting}{listing engine=minted, minted language=cpp, minted options={fontsize=\scriptsize}, listing only, colback=codebg, colframe=codebg, rounded corners}
// Код плагина

extern "C" IPlugin* create() {
  return new MyPlugin();
}

extern "C" void destroy(IPlugin* plugin) {
  delete plugin;
}
\end{tcblisting}
\end{itemize}
\end{frame}

\begin{frame}[fragile]
\frametitle{Плагины: shared\_ptr}
\begin{itemize}
\usemintedstyle{tango}
\item Удобно выделять объекты внутри плагина через \mintinline{cpp}|std::make_shared|: он запомнит \textit{deleter} при создании объекта, и подцепит функцию удаления из самого плагина
\pause
\item \textbf{NB:} Сам интерфейс плагина (\mintinline{cpp}|IPlugin|) так не создать: нетривиальные C++-объекты нельзя передавать/возвращать через \mintinline{cpp}|extern "C"| функции
\pause
\usemintedstyle{lightbulb}
\begin{tcblisting}{listing engine=minted, minted language=cpp, minted options={fontsize=\scriptsize}, listing only, colback=codebg, colframe=codebg, rounded corners}
// Код плагина
class MyPlugin : public IPlugin {
public:
  std::shared_ptr<Object> makeObject() override {
    return std::make_shared<Object>();
  }
};
\end{tcblisting}
\end{itemize}
\end{frame}

\begin{frame}[fragile]
\frametitle{Плагины: интерфейс через классы}
\begin{itemize}
\usemintedstyle{tango}
\item Экспортируем минимум функций
\pause
\item Все остальные функции плагина работают через понятный компилятору C++ интерфейс \begin{math}\Longrightarrow\end{math} сложнее ошибиться в использовании плагина
\pause
\item Фактически, мы заменяем импортирование неизвестных функций на механизм виртуальных функций C++
\end{itemize}
\end{frame}

\begin{frame}[fragile]
\frametitle{Boost}
\begin{itemize}
\usemintedstyle{tango}
\item Boost -- самый крупный набор библиотек общего назначения для С++
\pause
\item Обычно содержит код очень высокого качества, сравнимый со стандартной библиотекой, и работающий под большинством компиляторов
\pause
\item Многие части стандартной библиотеки пришли из Boost: \mintinline{cpp}|any, array, shared_ptr, chrono, ...|
\pause
\item Часто библиотеки из Boost ругают за излишнюю обобщённость (Boost.Geometry, Boost.GIL)
\end{itemize}
\end{frame}

\begin{frame}[fragile]
\frametitle{Boost}
\begin{itemize}
\usemintedstyle{tango}
\item Релизится 3 раза в год, первая версия вышла в 1999 году, последняя -- 3 апреля 2025 (версия 1.88.0)
\pause
\item \href{https://www.boost.org}{\nolinkurl{boost.org}}
\pause
\item Список библиотек: \href{https://www.boost.org/libraries/latest/grid}{\nolinkurl{boost.org/libraries/latest/grid}}
\end{itemize}
\end{frame}

\begin{frame}[fragile]
\frametitle{Сеть: IP, TCP, UDP}
\begin{itemize}
\usemintedstyle{tango}
\item IP (Internet Protocol) -- протокол сетевого уровня, связывающий разные компьютеры в сеть с помощью IP-адресов
\pause
\item TCP (Transmission Control Protocol) -- протокол транспортного уровня, обеспечивающий общение \underline{непрерывными потоками} данных, с гарантиями доставки
\pause
\item UDP (User Datagram Protocol) -- протокол транспортного уровня, обеспечивающий общение \underline{пакетами} данных, без каких-либо гарантий
\end{itemize}
\end{frame}

\begin{frame}[fragile]
\frametitle{Сеть: TCP vs UDP}
\begin{itemize}
\usemintedstyle{tango}
\item TCP гарантирует доставку данных и работает с ними в виде потока, это усложняет его реализацию и требует дополнительной синхронизации, приводящей к задержкам
\pause
\begin{itemize}
\item Весь интернет
\item Скачивание файлов
\item Общение между серверами
\end{itemize}
\pause
\item UDP ничего не гарантирует \begin{math}\Longrightarrow\end{math} он быстрее, но им сложнее пользоваться
\begin{itemize}
\item Мультиплеерные игры
\item Мультимедия (напр. стриминг видео)
\end{itemize}
\end{itemize}
\end{frame}

\begin{frame}[fragile]
\frametitle{Сеть: сокеты}
\begin{itemize}
\usemintedstyle{tango}
\item В большинстве ОС есть реализация POSIX сокетов -- интерфейса для работы с TCP/UDP
\pause
\item Сокеты -- особый тип файлового дексриптора, т.е. просто идентификатор типа \mintinline{cpp}|int|
\end{itemize}
\end{frame}

\begin{frame}[fragile]
\frametitle{Сеть: TCP сокеты}
\begin{itemize}
\usemintedstyle{tango}
\item Создать сокет и присоединиться к серверу по IP:
\usemintedstyle{lightbulb}
\begin{tcblisting}{listing engine=minted, minted language=cpp, minted options={fontsize=\scriptsize}, listing only, colback=codebg, colframe=codebg, rounded corners}
// SOCK_STREAM - TCP
// SOCK_DGRAM  - UDP
int my_socket = socket(AF_INET, SOCK_STREAM, 0);

sockaddr_in remote_address{};
remote_address.sin_family = AF_INET;
// Порт в network byte order
remote_address.sin_port = htons(80);
inet_pton(AF_INET, "23.215.0.138", &remote_address.sin_addr);

connect(my_socket, (sockaddr*)&remote_address,
  sizeof(remote_address));
\end{tcblisting}
\end{itemize}
\end{frame}

\begin{frame}[fragile]
\frametitle{Сеть: TCP сокеты}
\begin{itemize}
\item Отправка и получение данных -- как с файлами:
\usemintedstyle{lightbulb}
\begin{tcblisting}{listing engine=minted, minted language=cpp, minted options={fontsize=\scriptsize}, listing only, colback=codebg, colframe=codebg, rounded corners}
const char * message = ...;
write(my_socket, message, strlen(message));

char buffer[1024];
read(my_socket, buffer, sizeof(buffer));
\end{tcblisting}
\pause
\item Это \underline{блокирующие} операции!
\end{itemize}
\end{frame}

\begin{frame}[fragile]
\frametitle{Сеть: TCP сокеты}
\begin{itemize}
\item Привязать сокет к порту:
\usemintedstyle{lightbulb}
\begin{tcblisting}{listing engine=minted, minted language=cpp, minted options={fontsize=\scriptsize}, listing only, colback=codebg, colframe=codebg, rounded corners}
sockaddr_in server_address{};
server_address.sin_family = AF_INET;
server_address.sin_addr.s_addr = INADDR_ANY;
server_address.sin_port = htons(8080);

bind(my_socket, (struct sockaddr*)&server_address,
  sizeof(server_address));
listen(my_socket, 3);
\end{tcblisting}
\end{itemize}
\end{frame}

\begin{frame}[fragile]
\frametitle{Сеть: TCP сокеты}
\begin{itemize}
\item Принимать входящие соединения:
\usemintedstyle{lightbulb}
\begin{tcblisting}{listing engine=minted, minted language=cpp, minted options={fontsize=\scriptsize}, listing only, colback=codebg, colframe=codebg, rounded corners}
sockaddr_in client_address{};
socklen_t client_addrlen = sizeof(client_address);

while (true)
{
  int client_fd = accept(socket_fd,
    (struct sockaddr*)&client_address,
    &client_addrlen);

  // Общаемся через client_fd ...

  close(client_fd);
}
\end{tcblisting}
\pause
\item \textbf{NB:} Такой способ умеет работать только с одним клиентом одновременно
\end{itemize}
\end{frame}

\begin{frame}[fragile]
\frametitle{Сеть: много соединений}
\begin{itemize}
\item Часто программе нужно работать с несколькими соединениями параллельно
\pause
\item Банковское приложение: общается сразу с несколькими серверами
\pause
\item База данных: отвечает сразу на несколько параллельных запросов
\pause
\item Сервер онлайн-игры: отвечает сразу нескольким игрокам
\pause
\item И т.д.
\pause
\item Есть два основных способа это реализовать: потоки и асинхронные соединения
\end{itemize}
\end{frame}

\begin{frame}[fragile]
\frametitle{Сеть: потоки}
\begin{itemize}
\item Создаём по потоку на каждое соединение
\pause
\item Нужна синхронизация между потоками (мьютексы, атомики, и т.д.)
\pause
\item Переключение между потоками дорогое -- если соединений слишком много, будет много накладных расходов
\end{itemize}
\end{frame}

\begin{frame}[fragile]
\frametitle{Сеть: асинхронные соединения}
\begin{itemize}
\usemintedstyle{tango}
\item Асинхронная работа с сокетами требует поддержки со стороны ОС
\pause
\item Linux: \mintinline{cpp}|epoll|
\item Windows: \mintinline{cpp}|iocp| (I/O Completion Ports)
\item MacOS, FreeBSD: \mintinline{cpp}|kqueue|
\pause
\item Идея: отдаём ОС большой набор сокетов, и просим нас `разбудить', когда на какой-нибудь из них придут данные
\end{itemize}
\end{frame}

\begin{frame}[fragile]
\frametitle{Сеть: асинхронные соединения}
\begin{itemize}
\usemintedstyle{tango}
\item Потоки при этом использовать не обязательно (но можно)
\pause
\item Минимум накладных расходов, всё ожидание и синхронизацию за нас делает ОС
\end{itemize}
\end{frame}

\begin{frame}[fragile]
\frametitle{Сеть: асинхронные соединения}
\begin{itemize}
\usemintedstyle{tango}
\item Обычно это оборачивают в интерфейс, работающий с \textit{callback}'ами (например, лямбда-функциями)
\pause
\item Пример: отправь данные в сокет, а когда придёт ответ, вызови такую-то функцию:
\usemintedstyle{lightbulb}
\begin{tcblisting}{listing engine=minted, minted language=cpp, minted options={fontsize=\scriptsize}, listing only, colback=codebg, colframe=codebg, rounded corners}
socket.send(my_data, [&]{
  socket.receive(buffer, [&]{
    // handle data in buffer
  });
});
\end{tcblisting}
\pause
\item Такой код -- когда результат функции готов не сразу, а когда-нибудь -- называют \textit{асинхронным}
\end{itemize}
\end{frame}

\begin{frame}[fragile]
\frametitle{Сеть: асинхронные соединения}
\begin{itemize}
\item Асинхронность -- стандартный сегодня способ работы с сетью
\pause
\item Вместо callback'ов удобнее использовать \textit{корутины} -- они позволяют писать асинхронный код так же, как императивный
\pause
\item Node JS, Python asyncio, Boost.Asio, C++20 coroutines, etc
\end{itemize}
\end{frame}

\begin{frame}[fragile]
\frametitle{Сеть: Asio}
\begin{itemize}
\item Asio -- библиотека для асинхронной работы с сетью в C++
\pause
\item Входит в Boost, но есть и standalone вариант (т.е. не зависящий от Boost'а)
\pause
\item На её основе разрабатывается Networking TS -- работа с сетью в стандартной библиотеке C++
\pause
\item Поддерживает корутины из C++20
\end{itemize}
\end{frame}

\begin{frame}[fragile]
\frametitle{Сеть: Asio}
\begin{itemize}
\item Пример создания сокета и отправки данных:
\usemintedstyle{lightbulb}
\begin{tcblisting}{listing engine=minted, minted language=cpp, minted options={fontsize=\scriptsize}, listing only, colback=codebg, colframe=codebg, rounded corners}
using boost::asio;
using boost::asio::ip::tcp;

asio::io_context io_context;

tcp::socket socket(io_context);
connect(socket, address);
write(socket, buffer(message, sizeof(message)));
read(socket, buffer(answer, sizeof(answer)));
\end{tcblisting}
\end{itemize}
\end{frame}

\end{document}
